\documentclass{article}
\usepackage[a4paper, margin=0mm]{geometry}
\usepackage[usenames, dvipsnames]{xcolor}
\usepackage{tikz}
\usetikzlibrary{positioning, calc}

\usepackage{fontspec}
\newfontfamily\mono{UbuntuMono-R.ttf}

\usepackage{listings}
\lstset
{
    language=C,
    basicstyle=\mono\color{White},
    keywordstyle=\color{RoyalBlue},
    identifierstyle=\color{Plum},
    stringstyle=\color{Tan},
    commentstyle=\color{OliveGreen},
    directivestyle=\color{Tan},
    showstringspaces=false
}

\newcommand{\rslisting}[4]{\rotatebox[origin=c]{#1}{\scalebox{#2}{\lstinputlisting[identifierstyle=\color{#3}]{./../../../#4}}}}

\begin{document}
    \begin{tikzpicture}[remember picture,overlay]
        \fill[fill=darkgray] (current page.north west) rectangle (current page.south east);
        \node () at (20,0) {\rslisting{15}{1}{Aquamarine}{src/flag_detection.c}};
        \node () at (11,16) {\rslisting{0}{0.5}{Aquamarine}{src/flag_detection.c}};
        
        \node () at (8,6) {\rslisting{-6}{1}{Plum}{src/music.c}};
        \node () at (0,-15) {\rslisting{10}{1}{Plum}{src/music.c}};
     
        \node () at (20,0) {\rslisting{15}{1}{Aquamarine}{src/flag_detection.c}};
        \node () at (11,0) {\rslisting{20}{1}{Plum}{src/music.c}};
        \node () at (11,0) {\rslisting{20}{1.5}{BurntOrange}{src/main.c}};
        \node () at (15,-15) {\rslisting{-20}{2}{Tan}{src/state.c}};
        \node () at (8,5) {\rslisting{0}{1.2}{Emerald}{src/obstacle_detection.c}};
        \node () at (20,-20) {\rslisting{-45}{1.2}{Emerald}{src/dance.c}};
    \end{tikzpicture}
    
    \vspace*{\fill}
        \begin{center}
            \begin{tikzpicture}[remember picture, overlay]
                \fill[color = white, opacity = 0.9, rounded corners = 2ex] (-10,5) rectangle (10,-5);
                \node (titre) at (0,0)
                    {
                        \begin{tabular}{c}
                            \hline
                            \\
                            \huge\it Les aventures incroyables de \\
                            \\
                            \Huge\bf Jean-Bot Baptiste I\\
                            \\
                            \huge\it Le raccourceur des circuits\\
                            \\
                            \hline
                        \end{tabular}
                    };
                \node[left]  at ($(titre.west)-(0.2,0)$) {\includegraphics[width=5cm]{../images/kawwai.pdf}};
                \node[right] at ($(titre.east)+(0.2,0)$) {\scalebox{-1}[1]{\includegraphics[width=5cm]{../images/kawwai.pdf}}};
                \node[above] at (titre.north) {\textbf{\large MICRO-315 Miniprojet}};
                \node[below] at (titre.south) {\textit{\large Un projet de robotique de Julian Donevsky et Olof Lissmats}};
            \end{tikzpicture}
        \end{center}
        \vspace{5cm}
        \vspace*{\fill}
\end{document}